\documentclass[../../main.tex]{subfiles}
\begin{document}
% SOURCE p=176; printed=161; object=ch05-baire-category-title; region=middle
\section{Baire 纲定理及其推论}\label{section:ch05-baire-category}
% SOURCE p=176; printed=161; object=ch05-baire-category-introduction; region=middle
本节给出一个关于完备度量空间的重要定理，并用它得到 Banach 空间之间线性映射的若干基本结果。
% SOURCE p=176; printed=161; object=ch05-baire-category-theorem; number=5.9; region=middle
\begin{theorem}[Baire 纲定理]\label{theorem:ch05-baire-category}
设 $X$ 为完备度量空间。
\begin{enumerate}
\item\label{item:ch05-baire-dense-intersection} 若 $\{U_n\}_{n=1}^{\infty}$ 是 $X$ 中的一列开稠密子集，则 $\bigcap_{n=1}^{\infty}U_n$ 在 $X$ 中稠密。
\item\label{item:ch05-baire-not-meager} $X$ 不能表示为可数个无处稠密集的并。
\end{enumerate}
\end{theorem}
% SOURCE p=176; printed=161; object=ch05-baire-category-proof; region=middle-lower
\begin{proof}
对第一问，需证每个非空开集 $W\subseteq X$ 都与 $\bigcap_{n=1}^{\infty}U_n$ 相交。由于 $U_1\cap W$ 开且非空，其中包含一个球 $B(r_0,x_0)$，并可假设 $0<r_0<1$。对 $n>0$，按如下方式归纳选取 $x_n\in X$、$r_n\in(0,\infty)$：若 $j<n$ 的 $x_j,r_j$ 已选好，注意到 $U_n\cap B(r_{n-1},x_{n-1})$ 开且非空，故可选 $x_n,r_n$，使 $0<r_n<2^{-n}$ 且
\begin{align*}
\overline{B(r_n,x_n)}\subseteq U_n\cap B(r_{n-1},x_{n-1}).
\end{align*}
若 $n,m\ge N$，则 $x_n,x_m\in B(r_N,x_N)$；因 $r_n\to0$，序列 $\{x_n\}$ 为 Cauchy 列。由 $X$ 完备，极限 $x=\lim_nx_n$ 存在。由于 $n\ge N$ 时 $x_n\in B(r_N,x_N)$，故对所有 $N$ 有
\begin{align*}
x\in\overline{B(r_N,x_N)}\subseteq U_N\cap B(r_0,x_0)\subseteq U_N\cap W,
\end{align*}
第一问得证。

至于第二问，若 $\{E_n\}$ 为 $X$ 中的一列无处稠密集，则 $\{(\overline E_n)^c\}$ 是一列开稠密集。由于 $\bigcap_n(\overline E_n)^c\ne\varnothing$，有 $\bigcup_nE_n\subseteq\bigcup_n\overline E_n\ne X$。
\end{proof}
% SUPPLEMENT: id=v14-ch05-baire-details; parent=ch05-baire-category-theorem; kind=proof-detail+source-correction+motivation+hypotheses; origin=independent-derivation
\begin{remark}
\textbf{非空条件的范围（重构新增）。}第二问须假设 $X\ne\varnothing$：空空间完备，却等于可数个空集之并；空集在空空间内无处稠密。第一问对空空间仍成立，因为空集在空空间中稠密。原证明从“交稠密”推出“交非空”时，正是使用了底空间非空。下文关于“空间自身非第一纲”的结论采用此非空约定。
\end{remark}
\begin{proof}
\textbf{嵌套球与稠密性的细节补写（重构新增）。}任意非空开集 $V$ 中先取 $z\in V$ 和 $s>0$，使 $B(s,z)\subseteq V$；再取 $0<t<s/2$，则 $\overline{B(t,z)}\subseteq\{y:d(y,z)\le t\}\subseteq B(s,z)$。同时可把 $t$ 缩至小于任意指定正常数。这证明原构造中每一步的球确实存在。

对 $m,n\ge N$，三角不等式给 $d(x_m,x_n)<2r_N<2^{1-N}$，所以确为 Cauchy 列。各后续球及其中心均留在先前球内，闭包保留极限，所以 $x$ 同时属于全部 $U_N$ 和 $W$。因为每个非空开集 $W$ 都被交集遇到，交集稠密。无处稠密的定义给 $\operatorname{int}\overline E_n=\varnothing$，等价于其补为开稠密集，故非空情形的第二问由第一问推出。
\end{proof}
\begin{remark}
\textbf{动机与逐条件说明（重构新增）。}稠密性让每一步还能进入目标开集，开放性留下一个闭包仍在其中的小球，半径趋零使选出的中心为 Cauchy 列，完备性使它在空间内部有极限。可数性使这些要求可依次排列。完备性不能一般删除：$\mathbb Q$ 在通常度量下是可数个单点之并，各单点在 $\mathbb Q$ 中无处稠密。这里排除的是拓扑意义上的“小集”覆盖，与测度为零不是同一种小性，见\cref{exercise:ch05-meager-full-measure}。
\end{remark}
% END SUPPLEMENT: id=v14-ch05-baire-details
% SOURCE p=176; printed=161; object=ch05-baire-topological-invariance; region=lower
Baire 纲定理的结论纯属拓扑性质，因此只要 $X$ 同胚于一个完备度量空间就够了。例如定理适用于 $X=(0,1)$；它在通常度量下并不完备，但同胚于 $\mathbb R$。
% SOURCE p=176; printed=161; object=ch05-category-definitions; region=bottom
定理的名字来自 Baire 关于集合的术语。
\begin{definition}[第一纲、第二纲、疏集与剩余集]\label{definition:ch05-category}
若 $X$ 为拓扑空间，$E\subseteq X$ 是可数个无处稠密集的并，就称它为\textbf{第一纲}集；否则称为\textbf{第二纲}集。“第一纲”的一种较现代且更具描述性的同义词是\textbf{疏集}。疏集的补集称为\textbf{剩余集}。
\end{definition}
因此 Baire 定理断言，每个完备度量空间在自身中都是第二纲的。
% SUPPLEMENT: id=v14-ch05-category-details; parent=ch05-category-definitions; kind=proof-detail+well-definedness+motivation+hypotheses; origin=independent-derivation
\begin{proof}
\textbf{拓扑不变性及常用推论补写（重构新增）。}同胚把闭包、内部、开稠密集以及可数并和交分别映成对应对象，故保留无处稠密、疏和剩余性质，Baire 的两个结论也随之保留。函数 $t\mapsto\tan(\pi(t-1/2))$ 给 $(0,1)$ 到 $\mathbb R$ 的同胚，其连续逆为 $s\mapsto1/2+\pi^{-1}\arctan s$，验证了原例。

疏集的子集仍疏：若 $A\subseteq\bigcup_nE_n$ 且各 $E_n$ 无处稠密，则 $A=\bigcup_n(A\cap E_n)$，而每个 $A\cap E_n$ 的闭包内部包含于 $\operatorname{int}\overline E_n=\varnothing$。可数个疏集之并仍疏，把双重可数指标重新排列即可。

若 $X$ 为完备度量空间而 $W$ 为其中非空开集，则 $W$ 不可能作为 $X$ 的子集是疏的；否则写 $W\subseteq\bigcup_nE_n$，各 $(\overline E_n)^c$ 开稠密，Baire 第一问使其交遇到 $W$，与覆盖矛盾。因此剩余集在 $X$ 中稠密：若其补疏而它漏掉某个非空开集，该开集就会疏。特别，当 $X\ne\varnothing$ 时剩余集非空。
\end{proof}
\begin{remark}
\textbf{动机与条件说明（重构新增）。}这里的“纲”及疏性都相对于指定的环境 $X$ 判断；集合在自身中和在较大空间中可能有不同的疏性。剩余集描述可数种无处稠密障碍之外的点。它在完备度量空间中稠密，并不表示所有拓扑空间中的剩余集都非空。原文“每个完备度量空间为第二纲”仍须采用前述非空约定。
\end{remark}
% END SUPPLEMENT: id=v14-ch05-category-details
% SOURCE p=177; printed=162; object=ch05-baire-existence-applications; region=top
Baire 纲定理常用来证明存在性：在适当的完备度量空间里，只须证明不具有某性质的对象所成集合是疏集，就能证明具有该性质的对象存在。例如可由此证明处处不可微的连续函数存在，见\cref{exercise:ch05-nowhere-differentiable}。

现在转向 Baire 纲定理在线性映射理论中的应用。先介绍一些术语。
% SOURCE p=177; printed=162; object=ch05-open-mapping-definition; region=upper
\begin{definition}[开映射]\label{definition:ch05-open-map}
若 $X,Y$ 为拓扑空间，映射 $f:X\to Y$ 称为\textbf{开的}，若 $X$ 中每个开集 $U$ 的像 $f(U)$ 都在 $Y$ 中开。
\end{definition}
若 $X,Y$ 为度量空间，这等价于要求：每个以 $x\in X$ 为中心的球 $B$ 的像 $f(B)$，包含一个以 $f(x)$ 为中心的球。更特殊地，若 $X,Y$ 为赋范线性空间且 $f$ 线性，由于 $f$ 与平移和伸缩相容，$f$ 开当且仅当 $X$ 中以0为中心的单位球 $B$ 的像 $f(B)$ 包含 $Y$ 中一个以0为中心的球。
% SUPPLEMENT: id=v14-ch05-open-map-criterion; parent=ch05-open-mapping-definition; kind=proof-detail+motivation+hypotheses; origin=independent-derivation
\begin{proof}
\textbf{两个球判别的补写（重构新增）。}若 $f$ 开，则开球的像是包含中心像的开集，因而包含一个小开球。反过来，若球条件成立，给开集 $U$ 中每个 $x$ 取 $B(x,r_x)\subseteq U$，则 $f(U)$ 中的每个点 $f(x)$ 都有小球包含于 $f(B(x,r_x))\subseteq f(U)$，所以 $f(U)$ 开。

在线性情形，$B(x,r)=x+rB(0,1)$，故 $f(B(x,r))=f(x)+rf(B(0,1))$。若后者单位球像包含 $B(0,s)$，则 $f(B(x,r))$ 包含 $B(f(x),rs)$，应用前一判别即得开性。反方向取中心0、半径1即可。
\end{proof}
\begin{remark}
\textbf{动机与条件（重构新增）。}开性表示输入附近能覆盖输出附近，而连续性控制的是开集的原像；两者定义中方向不同。线性性把一个零点单位球的覆盖放大到所有点和尺度；无需预先假设 $f$ 连续。球始终指正半径的开球，且像内的球是在整个目标空间中取的。
\end{remark}
% END SUPPLEMENT: id=v14-ch05-open-map-criterion
% SOURCE p=177; printed=162; object=ch05-open-mapping-theorem; number=5.10; region=middle
\begin{theorem}[开映射定理]\label{theorem:ch05-open-mapping}
设 $\mathcal X,\mathcal Y$ 为 Banach 空间。若 $T\in L(\mathcal X,\mathcal Y)$ 满射，则 $T$ 为开映射。
\end{theorem}
% SOURCE p=177; printed=162; object=ch05-open-mapping-proof; region=middle-lower
\begin{proof}
用 $B_r$ 表示 $\mathcal X$ 中以0为中心、半径为 $r$ 的开球。由前面的说明，只须证明 $T(B_1)$ 包含 $\mathcal Y$ 中的一个零点球。由于 $\mathcal X=\bigcup_{n=1}^{\infty}B_n$ 且 $T$ 满射，有 $\mathcal Y=\bigcup_{n=1}^{\infty}T(B_n)$。但 $\mathcal Y$ 完备，且 $y\mapsto ny$ 是把 $T(B_1)$ 映到 $T(B_n)$ 的同胚，故由 Baire 定理，$T(B_1)$ 不能无处稠密。于是存在 $y_0\in\mathcal Y$、$r>0$，使 $B(4r,y_0)\subseteq\overline{T(B_1)}$。选 $y_1=Tx_1\in T(B_1)$，使 $\|y_1-y_0\|<2r$；则 $B(2r,y_1)\subseteq B(4r,y_0)\subseteq\overline{T(B_1)}$。所以若 $\|y\|<2r$，则
\begin{align*}
y=-Tx_1+(y+y_1)\in\overline{T(-x_1+B_1)}\subseteq\overline{T(B_2)}.
\end{align*}
两边除以2，得到存在 $r>0$，使 $\|y\|<r$ 时 $y\in\overline{T(B_1)}$。若能将这个闭包换为 $T(B_1)$，同时允许缩小 $r$，证明就完成了。下面实现这一点。

由于 $T$ 与伸缩相容，$\|y\|<r2^{-n}$ 时有 $y\in\overline{T(B_{2^{-n}})}$。设 $\|y\|<r/2$，可找 $x_1\in B_{1/2}$，使 $\|y-Tx_1\|<r/4$。依次归纳，可找 $x_n\in B_{2^{-n}}$，使
\begin{align*}
\left\|y-\sum_{j=1}^nTx_j\right\|<r2^{-n-1}.
\end{align*}
因为 $\mathcal X$ 完备，由\cref{theorem:ch05-banach-series-criterion}，级数 $\sum_{n=1}^{\infty}x_n$ 收敛，设其和为 $x$。于是 $\|x\|<\sum_{n=1}^{\infty}2^{-n}=1$，且 $y=Tx$。换言之，$T(B_1)$ 包含所有满足 $\|y\|<r/2$ 的 $y$，证毕。
\end{proof}
% SUPPLEMENT: id=v14-ch05-open-mapping-details; parent=ch05-open-mapping-theorem; kind=proof-detail+motivation+hypotheses; origin=independent-derivation
\begin{proof}
\textbf{闭包移位与逐次修正补写（重构新增）。}若 $T(B_1)$ 无处稠密，各正整数倍的像也无处稠密，满射给出的可数覆盖将与非空完备空间 $\mathcal Y$ 的 Baire 性矛盾。因此其闭包有非空内部，可在该内部选原证明所需的球。由 $y_0\in\overline{T(B_1)}$ 可选 $y_1$；平移是同胚，故
\begin{align*}
-Tx_1+\overline{T(B_1)}=\overline{-Tx_1+T(B_1)}=\overline{T(-x_1+B_1)}.
\end{align*}
而 $\|x_1\|<1$ 给 $-x_1+B_1\subseteq B_2$。正数倍乘也与闭包相容，所以除以2和后续的缩放均合法。

令初始残差 $e_0=y$。若 $\|e_{n-1}\|<r2^{-n}$，它在 $\overline{T(B_{2^{-n}})}$ 中，故可选 $x_n\in B_{2^{-n}}$，使 $e_n=e_{n-1}-Tx_n$ 满足 $\|e_n\|<r2^{-n-1}$。这是全部归纳步骤。又
\begin{align*}
\sum_{n=1}^{\infty}\|x_n\|\le\|x_1\|+\sum_{n=2}^{\infty}2^{-n}<\tfrac12+\tfrac12=1,
\end{align*}
所以极限 $x$ 真正在开单位球内，不只是闭单位球内。$T$ 有界给 $T(\sum_{j=1}^nx_j)\to Tx$，残差界同时使这个序列趋于 $y$，故 $Tx=y$。
\end{proof}
\begin{remark}
\textbf{动机与逐条件说明（重构新增）。}先用 $\mathcal Y$ 的完备性和满射的可数覆盖，找到单位球像闭包中的一个小球；再不断修正剩余误差，用 $\mathcal X$ 的完备性把近似原像相加成真正原像。线性性负责平移、缩放和叠加，有界性负责把无穷和传过 $T$。满射不能直接删除：例如 $\mathbb R\to\mathbb R^2$ 的坐标轴嵌入有界线性，单位区间的像在平面中不开。任一空间不完备时也可能失败，后面习题给具体例子。
\end{remark}
% END SUPPLEMENT: id=v14-ch05-open-mapping-details
% SOURCE p=177; printed=162; object=ch05-bounded-inverse; number=5.11; region=lower
\begin{corollary}\label{corollary:ch05-bounded-inverse}
若 $\mathcal X,\mathcal Y$ 为 Banach 空间且 $T\in L(\mathcal X,\mathcal Y)$ 双射，则 $T$ 为同构，即 $T^{-1}\in L(\mathcal Y,\mathcal X)$。
\end{corollary}
\begin{proof}
若 $T$ 双射，则 $T^{-1}$ 的连续性等价于 $T$ 的开性。
\end{proof}
% SUPPLEMENT: id=v14-ch05-bounded-inverse-details; parent=ch05-bounded-inverse; kind=proof-detail+motivation+hypotheses; origin=independent-derivation
\begin{proof}
\textbf{逆的线性与连续补写（重构新增）。}双射给逆映射存在；由 $T(ax+by)=aTx+bTy$ 和原像唯一性，逆映射线性。对 $\mathcal X$ 中开集 $U$，有 $(T^{-1})^{-1}(U)=T(U)$，开映射定理使它在 $\mathcal Y$ 中开，所以逆连续。由赋范空间上线性映射连续与有界等价，逆属于 $L(\mathcal Y,\mathcal X)$。
\end{proof}
\begin{remark}
\textbf{动机与条件说明（重构新增）。}在两个 Banach 空间间，原方向有界的线性双射自动在反方向也受控制。双射中的满射使开映射定理适用，单射使逆唯一；两端完备性继承自开映射定理。结论只保证一个有限逆界，并不声称原映射等距。
\end{remark}
% END SUPPLEMENT: id=v14-ch05-bounded-inverse-details
% SOURCE p=177; printed=162; object=ch05-closed-graph-definition; region=bottom
为得到下一个结果，再引入一些术语。
\begin{definition}[图像与闭线性映射]\label{definition:ch05-closed-graph}
若 $\mathcal X,\mathcal Y$ 为赋范向量空间，$T:\mathcal X\to\mathcal Y$ 为线性映射，定义 $T$ 的\textbf{图像}为
\begin{align*}
\Gamma(T)=\{(x,y)\in\mathcal X\times\mathcal Y:y=Tx\}.
\end{align*}
% SOURCE p=178; printed=163; object=ch05-closed-graph-definition-continuation; region=top
这是 $\mathcal X\times\mathcal Y$ 的一个子空间。（严格从集合论观点看，$T$ 与 $\Gamma(T)$ 本来相同；两者的区别只是思考方式不同。）若 $\Gamma(T)$ 为乘积空间的闭子空间，就称 $T$ 为\textbf{闭的}。
\end{definition}
显然，连续的 $T$ 总是闭的；若两空间完备，逆命题也成立。
% SUPPLEMENT: id=v14-ch05-closed-graph-definition-details; parent=ch05-closed-graph-definition; kind=proof-detail+well-definedness+motivation+hypotheses; origin=independent-derivation
\begin{proof}
\textbf{子空间及连续蕴含闭的补写（重构新增）。}图像含 $(0,0)$，且 $(x,Tx)$、$(y,Ty)$ 的线性组合为 $(ax+by,T(ax+by))$，仍在图像中。若 $T$ 连续且图像中的序列 $(x_n,Tx_n)$ 在乘积范数中趋于 $(x,y)$，则 $x_n\to x$、$Tx_n\to y$；连续性又给 $Tx_n\to Tx$，由极限唯一性得 $y=Tx$。所以图像序列闭，而乘积为度量空间，故图像闭。
\end{proof}
\begin{remark}
\textbf{动机与条件说明（重构新增）。}图像同时记录输入和输出；闭性保证一对收敛的输入、输出极限仍满足同一映射关系。乘积范数可取最大范数，也可取前述等价乘积范数。这里“闭映射”专指算子图像闭，不是拓扑学中“把闭集映为闭集”的映射性质。定义本身不要求 $T$ 有界。
\end{remark}
% END SUPPLEMENT: id=v14-ch05-closed-graph-definition-details
% SOURCE p=178; printed=163; object=ch05-closed-graph-theorem; number=5.12; region=upper
\begin{theorem}[闭图像定理]\label{theorem:ch05-closed-graph}
若 $\mathcal X,\mathcal Y$ 为 Banach 空间，$T:\mathcal X\to\mathcal Y$ 是闭的线性映射，则 $T$ 有界。
\end{theorem}
% SOURCE p=178; printed=163; object=ch05-closed-graph-proof; region=upper-middle
\begin{proof}
令 $\pi_1,\pi_2$ 为 $\Gamma(T)$ 到 $\mathcal X,\mathcal Y$ 的投影，即 $\pi_1(x,Tx)=x$、$\pi_2(x,Tx)=Tx$。显然 $\pi_1\in L(\Gamma(T),\mathcal X)$、$\pi_2\in L(\Gamma(T),\mathcal Y)$。由于 $\mathcal X,\mathcal Y$ 完备，乘积也完备；又因 $T$ 闭，$\Gamma(T)$ 也完备。$\pi_1$ 为从 $\Gamma(T)$ 到 $\mathcal X$ 的双射，由\cref{corollary:ch05-bounded-inverse}，$\pi_1^{-1}$ 有界。因此 $T=\pi_2\circ\pi_1^{-1}$ 有界。
\end{proof}
% SUPPLEMENT: id=v14-ch05-closed-graph-details; parent=ch05-closed-graph-theorem; kind=proof-detail+motivation+hypotheses; origin=independent-derivation
\begin{proof}
\textbf{投影和完备性的细节补写（重构新增）。}取乘积最大范数，两个坐标投影的范数都至多1；限制到图像仍有相同界。乘积中的 Cauchy 列给两个坐标各为 Cauchy 列，由两空间完备而有坐标极限；最大范数使坐标收敛等价于乘积收敛，所以乘积完备。图像闭，因此图像中的 Cauchy 列极限仍在图像内。

每个 $x\in\mathcal X$ 有且只有一个图像点 $(x,Tx)$，故 $\pi_1$ 确为双射。其定义域 $\Gamma(T)$ 及目标 $\mathcal X$ 都是 Banach 空间，满足有界逆推论的全部条件。复合界给 $\|T\|\le\|\pi_2\|\|\pi_1^{-1}\|<\infty$。
\end{proof}
\begin{remark}
\textbf{动机与逐条件说明（重构新增）。}把难以直接控制的 $T$ 换成图像上显然有界的两个投影，再借有界逆定理恢复输入到图像的控制。两端完备性保证乘积完备，图像闭性保证这个中间空间完备，线性性使图像及投影成为线性结构。$T$ 在整个 $\mathcal X$ 上有定义，使第一投影满射；只在真子空间上定义的算子不能直接套用这一定理。
\end{remark}
% END SUPPLEMENT: id=v14-ch05-closed-graph-details
% SOURCE p=178; printed=163; object=ch05-closed-graph-interpretation; region=middle
线性映射 $T:\mathcal X\to\mathcal Y$ 连续，意味着 $x_n\to x$ 时 $Tx_n\to Tx$；闭性则意味着若 $x_n\to x$ 且 $Tx_n\to y$，就有 $y=Tx$。所以闭图像定理的意义在于：验证 $x_n\to x$ 能推出 $Tx_n\to Tx$ 时，可以先假设 $Tx_n$ 已经收敛到某个值，然后只需证明极限是正确的值。这往往能省去许多麻烦。

两空间的完备性在开映射定理以及随后的闭图像定理中起关键作用。事实上，只要其中一个空间不完备，这两个定理的结论都可能失败，见\cref{exercise:ch05-weighted-incomplete,exercise:ch05-derivative-incomplete,exercise:ch05-unbounded-graph-space}。

本节最后一个结果具有几乎神奇的力量，允许在某些情形下从逐点估计推出一致估计。
% SOURCE p=178; printed=163; object=ch05-uniform-boundedness; number=5.13; region=middle-lower
\begin{theorem}[一致有界原理]\label{theorem:ch05-uniform-boundedness}
设 $\mathcal X,\mathcal Y$ 为赋范向量空间，$\mathcal A\subseteq L(\mathcal X,\mathcal Y)$。
\begin{enumerate}
\item\label{item:ch05-uniform-boundedness-nonmeager} 若在 $\mathcal X$ 的某个非疏子集中的每个 $x$ 都有 $\sup_{T\in\mathcal A}\|Tx\|<\infty$，则 $\sup_{T\in\mathcal A}\|T\|<\infty$。
\item\label{item:ch05-uniform-boundedness-banach} 若 $\mathcal X$ 为 Banach 空间，且每个 $x\in\mathcal X$ 都有 $\sup_{T\in\mathcal A}\|Tx\|<\infty$，则 $\sup_{T\in\mathcal A}\|T\|<\infty$。
\end{enumerate}
\end{theorem}
% SOURCE p=178; printed=163; object=ch05-uniform-boundedness-proof; region=lower
\begin{proof}
令
\begin{align*}
E_n=\{x\in\mathcal X:\sup_{T\in\mathcal A}\|Tx\|\le n\}
=\bigcap_{T\in\mathcal A}\{x\in\mathcal X:\|Tx\|\le n\}.
\end{align*}
则各 $E_n$ 闭；在第一问的假设下，某个 $E_n$ 必包含非平凡闭球 $\overline{B(r,x_0)}$。于是 $E_{2n}\supseteq\overline{B(r,0)}$，因为当 $\|x\|\le r$ 时，$x+x_0\in E_n$，故
\begin{align*}
\|Tx\|\le\|T(x+x_0)\|+\|Tx_0\|\le2n.
\end{align*}
换言之，只要 $T\in\mathcal A$ 且 $\|x\|\le r$，就有 $\|Tx\|\le2n$，因此 $\sup_{T\in\mathcal A}\|T\|\le2n/r$。第一问得证；第二问由 Baire 纲定理推出。
\end{proof}
% SUPPLEMENT: id=v14-ch05-uniform-boundedness-details; parent=ch05-uniform-boundedness; kind=proof-detail+motivation+hypotheses; origin=independent-derivation
\begin{proof}
\textbf{非疏集到统一球界的补写（重构新增）。}先假设 $\mathcal A$ 非空。每个 $T$ 连续且范数连续，故 $\{x:\|Tx\|\le n\}$ 闭，任意交 $E_n$ 闭。设逐点有界的非疏集为 $S$，则 $S\subseteq\bigcup_{n=1}^\infty E_n$。若每个 $E_n$ 内部都空，它们作为闭集就无处稠密，进而 $S$ 疏，矛盾。所以某个 $E_n$ 包含开球 $B(x_0,s)$。取 $0<r<s$，则 $\{x:\|x-x_0\|\le r\}\subseteq E_n$；这里“非平凡”只需理解为半径为正，零空间也允许。

原估计对所有 $T$ 使用同一个 $n,r$。对每个 $\|z\|\le1$ 令 $x=rz$，则 $r\|Tz\|\le2n$，对闭单位球取上确界即给算子范数统一界。第二问中，向量空间至少含零向量，因而非空；完备性和 Baire 定理给 $\mathcal X$ 在自身中非疏，取 $S=\mathcal X$ 应用第一问。若 $\mathcal A=\varnothing$，统一有界的全称断言真空成立；把非负数量在空族上的上确界约定为0，也使显示公式成立。
\end{proof}
\begin{remark}
\textbf{动机与逐条件说明（重构新增）。}逐点界原本可随 $x$ 变化，非疏性迫使某个共同整数界在整球上成立，再用线性平移和缩放扩展到所有单位向量。单个算子的有界性用于 $E_n$ 闭；算子族无需可数，因为闭集可以任意交，只有整数界 $n$ 需要可数。第一问不要求完备，第二问用定义域完备性保证足够多的测试点；目标 $\mathcal Y$ 不必完备，因为证明不在其中取新的极限。

定义域不完备时逐点有界可能不足：在有限支集序列空间 $c_{00}$ 上取上确界范数，令 $T_nx=nx_n$。每个 $T_n$ 有界且范数为 $n$；对固定有限支集的 $x$，只有有限多个 $T_nx$ 非零，故逐点有界，而范数无统一界。此空间确不完备，因为 $\sum_{j=1}^N2^{-j}e_j$ 是 Cauchy 列，其坐标极限有无限支集。
\end{remark}
% END SUPPLEMENT: id=v14-ch05-uniform-boundedness-details
\subfile{section-03-exercises.tex}
\end{document}
